% Chapter 11, generated by make_chapters.py from papers/Ch11_Bio_17_Inversions_Fibre/anccft30.tex.
% Source paper: Bio 17
\chapter[Inversions at inverted repeats]{Inversions at inverted repeats and the double-stranded fibre}\label{chap:ch11}
\begin{chaptersummary}
Two circular sequences with the same $k$-mer spectrum differ by transpositions of segments bounded by
repeated words: this is the theorem of Ukkonen and Pevzner. Double-stranded data do not see the
spectrum of one strand, only the sum of the spectra of both strands. The corresponding set of sequences
is the double-stranded fibre $\DS(S,k)$ counted in Chapters~\ref{chap:ch07}, \ref{chap:ch08} and~\ref{chap:ch09}. Chapter~\ref{chap:ch07} identified the move that
the second strand adds, the \emph{inversion flip}. Write $T=P\,wZ'\rc(w)\,R$ with $w$ a $k$-mer. Replacing
the segment $wZ'\rc(w)$ by its reverse complement preserves the double-stranded spectrum. For
$\phi$X174 at $k=11,12$ the flips alone connect the fibre. But flips alone cannot connect it when the
genome has no inverted repeat. We add the lattice picture: a flip moves the point of Chapter~\ref{chap:ch07} by
$-2(1-\rho)z$, where $z$ is a path from a vertex to its complement. Three results follow.
(i)~Transpositions and flips, together with global reverse complementation, generate the whole fibre
in every case we can close. The cases are twenty-four random sequences, checked against exhaustive
enumeration and against Chapter~\ref{chap:ch07}'s lattice sum, and $\phi$X174 at $k=12$ and $k=11$, where the closure of
the genome is exactly $\#\DS=10$ and $864$. Transpositions alone reach $2$ and $33$ sequences there,
and adding reverse complementation reaches $4$ and $66$. Adding transpositions removes Chapter~\ref{chap:ch07}'s
obstruction, and we state the general case as a conjecture.
(ii)~The \emph{inversion lattice} is always all of $\Lm$. The cycle part $(1-\rho)Z_1$ has index
\[
[\Lm:(1-\rho)Z_1]=\begin{cases}2&\text{if $D_k$ has no palindromic vertex and no self-complementary arc,}\\
1&\text{otherwise,}\end{cases}
\]
by a Tate-cohomology computation, confirmed on $600$ random covers. When the index is $2$, the missing
class is exactly an inversion vector.
(iii)~$\phi$X174 at $k=10$ carries $31$ inversions flanked by inverted repeats of length $10$. Random
walks of $20\,000$ moves stay inside the fibre of $31\,925\,753\,246\,212$ reconstructions and visit
about $5\,000$ of its lattice points. The inversion-group programme proposed for this chapter (Artin
groups, Garside normal forms) is set aside, and we say why.
\end{chaptersummary}

\section{The question}\label{ch11:sec:question}

Let $S$ be a circular DNA sequence and $k$ a word length. In the single-stranded world, the sequences
with the same $(k+1)$-mer spectrum as $S$ are the Eulerian circuits of one graph. Any two of them are
related by \emph{transpositions}: exchanges of segments that start and end at repeated $k$-mers
\cite{Ukk92,Pev95}. A sequencer that cannot tell the strands apart sees only
\[
\spec(T)+\spec(\rc T)=\spec(S)+\spec(\rc S),
\]
and the set $\DS(S,k)$ of circular $T$ satisfying this is larger. Chapter~\ref{chap:ch07} wrote its size as a sum over the
points of a lattice $\Lm$ in a box, Chapter~\ref{chap:ch08} counted the box, and Chapter~\ref{chap:ch09} computed
$\#\DS(\phi\mathrm X174,10)=31\,925\,753\,246\,212$. A transposition never
leaves the lattice point of $T$. So the question of this chapter is \emph{which moves carry one lattice point
to another}, and whether together with transpositions they connect the fibre.

The move itself is the \emph{inversion flip} of Definition~\ref{ch07:def:flip}: an inversion between two copies of an
inverted repeat. It is an operation every geneticist knows. In the human genome it is the classical
mechanism of recurrent inversions, via non-allelic homologous recombination between inverted low-copy
repeats \cite{SL02}, and here it is the combinatorial content of the second strand. Chapter~\ref{chap:ch07} proved two
things about it. The flips preserve the double-stranded spectrum, and a genome without inverted repeats
admits none, so its flip graph is discrete. It also computed that for $\phi$X174 at $k=12$ and $k=11$ the
flips alone connect all $10$ and $864$ reconstructions. This chapter adds three things. It combines flips
with transpositions, which removes that obstruction and gives a general conjecture. It computes the
lattice spanned by flip vectors (the Tate dichotomy). And it walks the fibre at $k=10$.

\section{Setting}\label{ch11:sec:setting}

We keep the notation of Chapter~\ref{chap:ch07}. $D_k$ is the labelled double cover of
Definition~\ref{ch07:def:cover}, whose arcs are the \emph{contents}; $\rho$ is reverse
complementation on vertices ($v\mapsto\bar v$) and on contents, a content $a\colon t\to h$ going to
$\rho a\colon\bar h\to\bar t$. A reconstruction $T\in\DS(S,k)$ is an Eulerian circuit of the multigraph
with $c(a)$ copies of each content, where $c+\rho c=m$, and $\#\DS(S,k)$ is the lattice sum of
Theorem~\ref{ch07:thm:lattice} over the box points $f=2c-m$ of $\Lm$. Here $(\rho z)(a)=z(\rho a)$ on $1$-chains. A vertex is \emph{palindromic} if $\bar v=v$, and
a content is \emph{self-complementary} if $\rho a=a$.

\section{Inversions at inverted repeats}\label{ch11:sec:inv}

\begin{theorem}[flips preserve the double-stranded spectrum {Proposition~\ref{ch07:prop:flip}(a)}]\label{ch11:thm:inv}
Let $T$ be circular and $w$ a $k$-mer with an occurrence of $w$ and a later occurrence of $\rc(w)$, so
that, read from the first, $T=Z\,R$ with $Z=w\,Z'\,\rc(w)$. Let $T'=\rc(Z)\,R$. Then
$\spec_{k+1}(T')+\spec_{k+1}(\rc T')=\spec_{k+1}(T)+\spec_{k+1}(\rc T)$.
\end{theorem}

\begin{proof}
$\rc(Z)=w\,\rc(Z')\,\rc(w)$ begins with $w$ and ends with $\rc(w)$, as $Z$ does. A $(k+1)$-mer of $T$ that
straddles a junction between $Z$ and $R$ meets $Z$ in at most $k$ letters, and those letters lie inside
the flanking copy of $w$ or $\rc(w)$. So it is also a $(k+1)$-mer of $T'$. The $(k+1)$-mers inside $Z$ are
replaced by their reverse complements, so
$\spec(T')=\spec(T)-\spec(Z)+\rc\,\spec(Z)$. Applying $\rc$ and adding gives the claim.
\end{proof}

This is Proposition~\ref{ch07:prop:flip}(a), with the same proof; we repeat it only to fix
notation. The move of this section is a \emph{flip} in the sense of Definition~\ref{ch07:def:flip}.

On the double cover, the two copies $w$, $\rc(w)$ are vertices $v$, $\bar v$. The segment $Z$ is a path
$z$ of contents from $v$ to $\bar v$, and $\rho z$ is again a path from $v$ to $\bar v$, because $\rho$
reverses paths and exchanges their ends. The inversion replaces $z$ by $\rho z$:
\begin{equation}\label{ch11:eq:move}
c'=c-z+\rho z,\qquad f'=f-2(1-\rho)z .
\end{equation}
The vector $(1-\rho)z$ is a cycle, since $z$ and $\rho z$ have the same ends, and it is $\rho$-antisymmetric.
So it lies in $\Lm$. An inversion changes the strand of every content of $Z$. A transposition changes no
multiplicity, so $f'=f$.

Check (I) of \texttt{bio17\_verify.py} tests Theorem~\ref{ch11:thm:inv} directly on strings. All $156$
inversions at inverted repeats of the random test sequences preserve the double-stranded spectrum. So
do all $2$, $9$ and $31$ such inversions of $\phi$X174 itself at $k=12$, $11$ and $10$. As a control,
only $64$ of $960$ inversions of random, unflanked segments preserve it, and those are chance
coincidences at small $k$.

\section{Which moves generate the fibre}\label{ch11:sec:gen}

We represent a reconstruction as a cyclic sequence of contents and apply three families of moves:
\begin{itemize}
\item \textbf{T} (transpositions): at a vertex visited three times, exchange two of the returning
loops; or, at two vertices $u\ne w$ visited alternately $u\,w\,u\,w$, exchange the two $u\to w$
segments;
\item \textbf{I}: the inversions of Theorem~\ref{ch11:thm:inv}, at every pair of visits to $v$ and $\bar v$;
\item \textbf{R}: global reverse complementation.
\end{itemize}
The closure of $S$ under a family is computed by breadth-first search, and every state is checked to be a
valid reconstruction.

\paragraph{Three counts agree.} On twenty-four random sequences of length $16$--$33$ with $k=3$--$5$,
chosen so that at least two lattice points carry reconstructions, three quantities coincide in every
case. The first is the set of circular strings with the double-stranded spectrum of $S$, found by
exhaustive search. The second is Chapter~\ref{chap:ch07}'s lattice sum. The third is the closure of $S$ under T, I and R.
The fibres range from $4$ to $1\,872$ sequences (check~(C)).

\begin{table}[t]
\centering\small
\begin{tabular}{lcccc}
\toprule
& T & T $+$ R & T $+$ I & T $+$ I $+$ R\\
\midrule
random cases whose closure is all of $\DS$ (of $24$) & 0 & 2 & 23 & 24\\
$\phi$X174, $k=12$ ($\#\DS=10$) & 2 & 4 & 10 & 10\\
$\phi$X174, $k=11$ ($\#\DS=864$) & 33 & 66 & 864 & 864\\
\bottomrule
\end{tabular}
\caption{Closures of $S$ under the move families: T transpositions, R reverse complement, I inversions
at inverted repeats. For $\phi$X174 the entries are the sizes of the closures. At $k=12$ and $k=11$ the full
closure visits $8$ and $218$ lattice points, exactly the box points with connected support found by
Chapter~\ref{chap:ch07} and Chapter~\ref{chap:ch08}.}\label{ch11:tab:gen}
\end{table}

For $\phi$X174 these closures refine Theorem~\ref{ch07:thm:flips}, which showed that at $k=12$ and
$k=11$ the flips \emph{alone}, without transpositions and without reverse complementation, already
connect the fibre. What transpositions add is the general case:
Proposition~\ref{ch07:prop:flip}(d) shows that a genome without inverted repeats admits no flip at all.
Table~\ref{ch11:tab:gen} separates the contributions. Transpositions alone give the single-stranded fibre of
$S$: $\Nseq$ at its own lattice point, $2$ and $33$ for $\phi$X174. Reverse complementation doubles
this. The inversions are what carry the closure across the lattice. At $k=11$ they alone, with
transpositions, reach all $864$ reconstructions and all $218$ lattice points with connected support. In
one random case of twenty-four, T and I reach only part of the fibre and R is needed as well. By
Proposition~\ref{ch07:prop:flip}(c), reverse complementation is itself a flip only at a palindromic visit. We have not
checked whether the exceptional case lacks one.

\begin{conjecture}[double-stranded Ukkonen--Pevzner]\label{ch11:conj:gen}
For every circular $S$ and every $k$, $\DS(S,k)$ is connected under transpositions at repeated
$k$-mers, inversions at inverted repeats of length $k$, and global reverse complementation.
\end{conjecture}

Within one lattice point the conjecture is the classical theorem, and on the palindromic genomes of
Section~\ref{ch10:sec:real} it holds by Kotzig's theorem (Section~\ref{ch10:sec:by}). The open part is the passage
between lattice points: whether any two box points with connected support are joined by a chain of
inversions that are realizable, i.e.\ whose flanking copies $v$, $\bar v$ are actually visited by the
current circuit. Section~\ref{ch11:sec:lattice} shows that nothing is lost at the level of the lattice.

\section{The inversion lattice}\label{ch11:sec:lattice}

By \eqref{ch11:eq:move} the inversion vectors $(1-\rho)z$, for all paths $z$ from a vertex to its complement,
generate a sublattice $M\subseteq\Lm$. The differences of two such paths with the same ends are cycles,
so $M$ contains $(1-\rho)Z_1$. The question is whether $M$, or already $(1-\rho)Z_1$, is all of $\Lm$.
The group $\langle\rho\rangle\cong\Z/2$ acts, and the answer is a Tate-cohomology computation. Recall
that for $\Z/2$ acting on a module $A$ one has
$\Hh^0(A)=\ker(1-\rho)/\operatorname{im}(1+\rho)$ and $\Hh^1(A)=\ker(1+\rho)/\operatorname{im}(1-\rho)$
\cite{Bro82}. In particular $\Lm/(1-\rho)Z_1=\Hh^1(Z_1)$.

\begin{theorem}[the Tate dichotomy]\label{ch11:thm:tate}
Let $D_k$ be connected. Then
\[
\Lm/(1-\rho)Z_1\;\cong\;\begin{cases}\Z/2&\text{if $D_k$ has no palindromic vertex and no
self-complementary content,}\\ 0&\text{otherwise.}\end{cases}
\]
In every case $M=\Lm$. In the first case the non-trivial class is represented by $(1-\rho)p$ for any
path $p$ from a vertex $v$ to $\bar v$.
\end{theorem}

\begin{proof}
Let $C_1$ be the free abelian group on the contents and $C_0$ that on the vertices, with
$\partial a=h(a)-t(a)$. Because $\rho a$ runs from $\bar h$ to $\bar t$, one has
$\partial\rho=-\rho\,\partial$. So $\partial$ is equivariant when $C_0$ carries the twisted action
$\sigma=-\rho$. Let $B_0=\partial C_1$. Since $D_k$ is connected, $B_0$ is the group of $0$-chains of
total degree $0$. The sequence $0\to Z_1\to C_1\to B_0\to0$ is exact and equivariant, and its long exact
sequence in Tate cohomology reads
\[
\Hh^0(C_1)\longrightarrow\Hh^0(B_0)\xrightarrow{\ \delta\ }\Hh^1(Z_1)\longrightarrow\Hh^1(C_1).
\]
\emph{$C_1$.} The contents fall into free orbits $\{a,\rho a\}$, each a copy of $\Z[\Z/2]$ with vanishing
Tate cohomology, and self-complementary contents, each a copy of the trivial module $\Z$, with
$\Hh^0=\Z/2$ and $\Hh^1=0$. So $\Hh^1(C_1)=0$, and $\Hh^0(C_1)=(\Z/2)^{\#\text{self-compl.}}$, generated by
the classes of the self-complementary contents.

\emph{$B_0$.} $\ker(1-\sigma)$ consists of the $x$ with $x(\bar v)=-x(v)$. Such an $x$ vanishes at palindromic
vertices, and is spanned by the $e_v-e_{\bar v}$. Now $\operatorname{im}(1+\sigma)$ consists of the
$y-\rho y$ with $y\in B_0$. If $p$ is a palindromic vertex, then $e_v-e_{\bar v}=(1-\rho)(e_v-e_p)$ with
$e_v-e_p\in B_0$, so $\Hh^0(B_0)=0$. If there is none, any $y$ with $(1-\rho)y=e_v-e_{\bar v}$ has the form
$e_v+s$ with $\rho s=s$, and $s$ has even total degree, because its values pair up over orbits of size
two. So $y\notin B_0$. Hence $\Hh^0(B_0)=\Z/2$, generated by the common class of all $e_v-e_{\bar v}$;
they differ by $(1-\rho)(e_v-e_u)$ with $e_v-e_u\in B_0$.

\emph{The map.} A self-complementary content $a\colon t\to h$ has $h=\bar t$, so
$\partial a=e_{\bar t}-e_t$, the generator of $\Hh^0(B_0)$. So if a self-complementary content exists,
$\Hh^0(C_1)\to\Hh^0(B_0)$ is onto and $\Hh^1(Z_1)=0$. If none exists and there is no palindromic vertex,
$\Hh^1(Z_1)\cong\Hh^0(B_0)=\Z/2$. If there is a palindromic vertex, $\Hh^0(B_0)=0$ and again
$\Hh^1(Z_1)=0$.

\emph{The representative.} The connecting map sends the class of $x=e_{\bar v}-e_v$ to the class of
$z-\rho z$, for any $z\in C_1$ with $\partial z=x$: then $\partial(\rho z)=-\rho x=x$, so $z-\rho z\in Z_1$. A
path $p$ from $v$ to $\bar v$ is such a $z$. So $(1-\rho)p$ represents the non-trivial class, and $M=\Lm$.
\end{proof}

\begin{table}[t]
\centering\small
\begin{tabular}{lrcc}
\toprule
double cover & covers & $[\Lm:(1-\rho)Z_1]$ & $[\Lm:M]$\\
\midrule
no palindromic vertex, no self-complementary content & 60 & $2$ in all $60$ & $1$\\
self-complementary content, no palindromic vertex & 241 & $1$ & $1$\\
palindromic vertex, no self-complementary content & 299 & $1$ & $1$\\
\midrule
$\phi$X174, $k=10$ ($3$ palindromic vertices) & 1 & 1 & 1\\
$\phi$X174, $k=11$ ($1$ self-complementary content) & 1 & 1 & 1\\
$\phi$X174, $k=12$ ($1$ palindromic vertex) & 1 & 1 & 1\\
\bottomrule
\end{tabular}
\caption{Indices computed by Smith normal form on $600$ random connected double covers (lengths
$14$--$59$, $k=2$--$7$; $36$ disconnected covers skipped) and on $\phi$X174. The ranks of both lattices
always equal $\operatorname{rank}\Lm$: $48$, $12$ and $4$ for $\phi$X174.}\label{ch11:tab:tate}
\end{table}

Table~\ref{ch11:tab:tate} confirms the theorem exactly on every cover tested. The dichotomy has a transparent
meaning. A palindromic vertex, or a self-complementary content, is a place where the two strands
touch. Without one, flipping closed cycles, $(1-\rho)Z_1$, can only reach a sublattice of index $2$, and
the missing half needs an open path from a word to its complement: a genuine inversion.

\section{\texorpdfstring{$\phi$X174 at $k=10$}{phiX174 at k=10}}\label{ch11:sec:phix}

At $k=10$ the fibre has $31\,925\,753\,246\,212$ elements on $1\,240\,843\,950\times2$ box points
(Chapter~\ref{chap:ch09}). It cannot be closed, but it can be walked. The genome carries $31$ inversions at inverted
repeats of length $10$. The count on the cover is $32$. The extra one comes from a copy of $\rc(w)$ at
position $4482$ and a copy of $w$ at $4484$. Their two flanks overlap, and the string-level count
excludes that case. One inversion from $\phi$X174 reaches $23$ other lattice
points. Two random walks of $20\,000$ moves each, one without and one with global reverse
complementation, stay inside the fibre at every step: each state is checked to be a valid
reconstruction, and the final one also at the level of strings. They visit $4\,907$ and $5\,049$
distinct lattice points, and none of the sampled sequences repeats. That is consistent with
Conjecture~\ref{ch11:conj:gen} but does not test it: at this scale the question is not whether the walk can
move, but whether it can reach every one of the $454\,551\,605\times2$ points with connected support.

\section{Biological reading}\label{ch11:sec:bio}

Theorem~\ref{ch11:thm:inv} says exactly which rearrangements double-stranded $k$-mer data cannot see:
transpositions between direct repeats, and inversions between inverted repeats, each at least $k$ long.
These are the two classical products of non-allelic homologous recombination: deletions,
duplications and translocations between direct repeats, and inversions between inverted ones
\cite{SL02}. The first class preserves the content of one strand and the second only the content of
both. Any strand-agnostic encoder that depends on $k$-mer content alone is blind to both. That includes
a reverse-complement-invariant local DNA encoder of window $k+1$, in the sense of
Section~\ref{ch12:sec:dna}. A short-read assembler faces the same limit. The inversion lattice being
all of $\Lm$ means that, algebraically, these two mechanisms already account for every direction in
which the double-stranded fibre extends.

\section{On the proposed inversion-group programme}\label{ch11:sec:programme}

This chapter was specified with a different programme: a group $G_{\mathrm{inv}}$ generated by adjacent
exchanges and segment inversions, embedded in an Artin--Tits group, shown to be Garside, with the
minimum inversion distance read off a Garside normal form. We set it aside for three reasons.
\begin{itemize}
\item The group generated by signed permutations of $n$ oriented segments is the finite hyperoctahedral
group $B_n$, of order $2^nn!$. It is a Coxeter group, not an Artin group, and it does not embed in one
as the specification intended. Its word problem is trivial.
\item The minimum number of inversions between two signed orders, the reversal distance, is computed in
polynomial time by Hannenhalli and Pevzner \cite{HP99}. Word length in $B_n$ with respect to its
Coxeter generators is a different quantity, and it does not model inversions of arbitrary segments.
\item The object that is new, and that belongs to algebraic genomics, is not the group of all inversions
but the inversions that preserve the double-stranded spectrum. Those are the ones flanked by inverted
repeats, and they are what this chapter studies.
\end{itemize}

\section*{Ancillary file}
\texttt{bio17\_verify.py} (standard library; imports \texttt{ds\_fermion.py} of Chapter~\ref{chap:ch09}) runs checks
(I), (C), (G), (L), (H) and (W). These are the inversion theorem on strings; the three counts; the
closures; the inversion lattice on $\phi$X174 and random covers; the Tate dichotomy on $600$ random
covers; and the walks at $k=10$. The whole battery takes under two minutes.
